% Part: normal-modal-logic
% Chapter: filtrations
% Section: preliminaries

\documentclass[../../../include/open-logic-section]{subfiles}

\begin{document}

\olfileid[tr]{nml}{fil}{pre}

\olsection{Ön Bilgiler}

Süzmeler, modal mantık sistemlerimizin \emph{sonlu model özelliğine} sahip
olduğunu, yani bir modelde doğru (yanlış) olan her !!{formula}{}ün bir
\emph{sonlu} modelde de doğru (yanlış) olduğunu göstererek bu sistemlerin
karar verilebilirliğini kurmamızı sağlar. Süzmeler, alt formüller altında
kapalı !!{formula}s kümelerine göre tanımlanır.

\begin{defn}\ollabel{defn:modallyclosed}
  Bir $\Gamma$ !!{formula}s kümesi, $\Gamma$ içindeki !!a{formula}{}ün her
  alt formülünü içeriyorsa \emph{alt formüller altında kapalıdır}. Ayrıca
  $\Gamma$ alt formüller altında kapalıysa ve $!A\in\Gamma$ olması
  $\Box!A,\Diamond!A\in\Gamma$ olmasını da gerektiriyorsa $\Gamma$
  \emph{modal olarak kapalıdır}.
\end{defn}

Örneğin !!a{formula} $!A$ verildiğinde, onun bütün alt !!{formula}s{}inin
kümesi alt !!{formula}s altında kapalıdır. Bir modelin $!A$'nın alt
!!{formula}s{}i kümesinden süzülmesini tanımladığımızda istediğimiz özelliğe
sahip olur: $!A$'yı ancak ve ancak özgün model doğru (yanlış) yapıyorsa doğru
(yanlış) yapar.

$\mModel{M}$ modelinin $\Gamma$ üzerinden bir süzmesinin dünyalar kümesi,
aşağıdaki denklik bağıntısının bütün denklik sınıfları kümesi olarak
tanımlanır.

\begin{defn}
$\mModel{M}=\tuple{W,R,V}$ olsun ve $\Gamma$ alt !!{formula}s altında
kapalı olsun. $W$ üzerinde, $\Gamma$'dan tam olarak aynı !!{formula}s{}i doğru
yapan herhangi iki dünya arasında geçerli bir $\equiv$ bağıntısı tanımlayın:
\[
u \equiv v \quad \text{ancak ve ancak} \quad
\forall !A \in \Gamma : \mSat{M}{!A}[u] \Leftrightarrow \mSat{M}{!A}[v].
\]
Bir $w$ dünyasının $[w]_\equiv$ denklik sınıfı, kısaca $[w]$, $w$'ye
$\equiv$-denk bütün dünyaların kümesidir:
\[
[w] = \Setabs{v}{v \equiv w}.
\]
\end{defn}

\begin{prop}
  $\mModel{M}$ ve $\Gamma$ verildiğinde, yukarıda tanımlanan $\equiv$ bir
  denklik bağıntısıdır; yani yansımalı, simetrik ve geçişlidir.
\end{prop}

\begin{proof}
  $w$, $\Gamma$ içindeki tam olarak kendisiyle aynı !!{formula}s{}i doğru
  yaptığından $\equiv$ yansımalıdır. $u$, $\Gamma$ içindeki $v$ ile aynı
  !!{formula}s{}i doğru yapıyorsa bunun $v$ ile $u$ için de geçerli olması
  nedeniyle simetriktir. $u$, $\Gamma$ içindeki $v$ ile aynı; $v$ de $w$
  ile aynı !!{formula}s{}i doğru yapıyorsa $u$ ile $w$'nin de aynı
  !!{formula}s{}i doğru yapması nedeniyle geçişlidir.
\end{proof}

Her denklik bağıntısı gibi $\equiv$ de $W$'yi \emph{bölümlere}, yani
$W$'nin ikişer ikişer ayrık ve birlikte $W$'nin tamamını kapsayan alt
kümelerine ayırır. Her $w\in W$, $w\equiv w$ olduğundan bölümlerden birinin,
yani $[w]$'nin !!a{element}{}dır. Dolayısıyla $[w]$ bölümleri $W$'nin
tamamını kapsar. İkişer ikişer ayrıklardır; çünkü $u\in[w]$ ve $u\in[v]$
ise $u\equiv w$ ve $u\equiv v$; simetriklik ve geçişlilikten $w\equiv v$,
dolayısıyla $[w]=[v]$ çıkar.

\end{document}
